\documentclass[a4paper,11pt]{article}
\usepackage{exam-style-842-unit}
\examinehead{842 信息化概论 · 第07单元练习册}{练习一 参考答案与解析}
\begin{document}

\unitanswerscover{842 信息化概论 · 单元练习册}{第07单元参考答案}{练习一（人工智能基础）的答案与解析；每道判断题均给出改正后的正确表述，教材未展开处已注明来源。}

\answerpart{一、单项选择题（10 题，每题 3 分，共 30 分）}

\q A。\begin{analysis}教材 9.1（印刷页 149）指出：一般认为，智能是指人类在认识世界和改造世界的活动中，由脑力劳动表现出来的能力，即个体对客观事物进行合理分析、判断及有目的的行动和有效处理周围环境事宜的综合能力。B 描述的是计算机执行程序的能力，C 是生物本能，D 只涉及知识数量，都不是智能的定义。\end{analysis}

\q B。\begin{analysis}美国斯坦福大学人工智能研究中心的尼尔逊教授认为，人工智能是关于知识的学科，是怎样表示知识以及怎样获得知识并使用知识的科学（教材 9.2，印刷页 150）。A 图灵提出图灵机与图灵测试；C 麦卡锡在 1956 年达特茅斯会议上提议采用“人工智能”术语；D 费根鲍姆 1977 年提出知识工程。\end{analysis}

\q B。\begin{analysis}图灵测试由图灵于 1950 年在《计算机与智能》中提出：测试者通过文字终端分别与被测试的人和机器对话，若无法区分对方是人还是机器，则认为机器通过了测试、表现出智能。教材第 1 章 1.3.4（印刷页 11—12）只提到图灵“提出了一种用于判定机器是否具有智能的试验方法”，未展开具体做法（该处 OCR 有残缺），完整表述按通用教材补充；教材第 9 章 9.7.1（印刷页 162）记载 1950 年图灵出版《计算机与智能》。A、C、D 分别是运算速度、定理证明和存储容量，均不是图灵测试的内容。\end{analysis}

\q C。\begin{analysis}1956 年，麦卡锡等人在达特茅斯学院举行的一次会议上提议正式采用“人工智能（AI）”这一术语，标志着人工智能学科诞生（教材 9.3，印刷页 151）。A 是孕育期的人工神经元模型，B 是图灵的文章，D 是知识工程提出的会议。\end{analysis}

\q C。\begin{analysis}1977 年费根鲍姆提出知识工程后，各类专家系统得以发展，大量商品化专家系统和智能系统纷纷推出，这一阶段是知识应用期（教材 9.3，印刷页 151—152）。孕育期的标志是人工神经元模型；形成期以 1956 年达特茅斯会议为标志；综合集成期从 20 世纪 80 年代末开始。\end{analysis}

\q B。\begin{analysis}1977 年，费根鲍姆在第五届国际人工智能联合会议上正式提出了知识工程的概念（教材 9.3，印刷页 151）。A 麦卡锡与“人工智能”术语、C 布鲁克斯与行为模拟机器人、D 尼尔逊与“关于知识的学科”定义分别对应，不能混淆。\end{analysis}

\q A。\begin{analysis}教材 9.4.1（印刷页 152）指出，许多问题（包括智力问题和实际工程问题）的求解都可以归结为对某种图或空间的搜索问题，因此搜索技术就成为人工智能最基本的研究内容。B、C、D 都是人工智能的主要内容，但不是“最基本”的研究内容。\end{analysis}

\q C。\begin{analysis}计算智能以数值数据为基础，主要经数值计算、运用算法求解，主要内容包括神经计算、进化计算、免疫计算、粒群算法、蚁群算法、自然计算和人工生命等（教材 9.2.4，印刷页 151）。A 知识表示、B 谓词逻辑推理、D 专家系统都属于符号智能或知识工程范畴。\end{analysis}

\q A。\begin{analysis}教材 9.2.2（印刷页 150）指出，智能机器至少要具备三个要素：感觉要素、反应要素和思考要素。注意教材 9.6.15（印刷页 160—161）在讲智能机器人时表述为“感觉要素、运动要素、思考要素”，前后略有差异；本题按 9.2.2 作答。B、C、D 均非教材表述。\end{analysis}

\q D。\begin{analysis}三者是包含关系：人工智能是最大的学科门类，机器学习是人工智能的一个分支，深度学习是机器学习的一个分支，通常以多层人工神经网络为模型（教材 9.6.14 给出“深度学习是多层的人工神经网络和训练的方法”，教材未系统展开三者层次，按通用教材补充）。A 把人工智能等同于机器学习，B 把深度学习与神经网络等同，C 颠倒包含关系，均错。\end{analysis}

\answerpart{二、多项选择题（4 题，每题 4 分，共 16 分）}

\q ABC。\begin{analysis}教材 9.1（印刷页 149）把智能本质的共识概括为三点：（1）智能具有感知能力；（2）智能具有记忆和思维能力；（3）智能具有学习能力、自适应能力及行为能力。D 错误：情感不是教材所列智能本质的共识内容（教材 9.7.3 只是把“机器情感”列为未来可能的发展方向之一）。\end{analysis}

\q ABC。\begin{analysis}符号智能以符号形式的知识和信息为基础，主要内容包括知识获取、知识表示、知识组织与管理和知识运用等技术，构成知识工程及基于知识的智能系统（教材 9.2.4，印刷页 150—151）。D 错误：数值计算与优化搜索是计算智能的主要内容，不属于符号智能。\end{analysis}

\q ABC。\begin{analysis}A 正确：心理模拟符号推演从人脑的宏观心理层面入手（教材 9.5.1）；B 正确：生理模拟神经计算用人工神经网络作为信息和知识的载体（教材 9.5.2）；C 正确：行为模拟法的典型代表是 MIT 布鲁克斯教授研制的六足行走机器人（教材 9.5.3）。D 错误：群体模拟仿生计算模拟的是生物群落的群体智能行为，如遗传算法、蚁群算法、粒群算法等，不是单个神经元的工作机理。\end{analysis}

\q ABC。\begin{analysis}A 正确：1956 年达特茅斯学院会议正式确立了人工智能的研究领域（教材 9.3）；B 正确：1977 年费根鲍姆提出知识工程；C 正确：1969 年国际人工智能联合会议标志着人工智能得到国际认可（教材 9.3，印刷页 151）。D 错误：教材明确写道人工智能的发展“高潮和低谷不断交替出现”，如 1973 年美英政府停止向没有明确目标的人工智能研究项目拨款、20 世纪 80 年代末投资者重新撤回投资。\end{analysis}

\answerpart{三、判断题（6 题，每题 2 分，共 12 分）}

\q ×。改正：智能包括学习能力，学习能力是指通过指导、实践等过程来丰富自身知识和技巧的能力，后天学习和实践是智能的重要来源。\begin{analysis}教材 9.1（印刷页 149）把学习能力、自适应能力及行为能力列为智能的第三点共识，可见智能并非与学习和实践无关的先天遗传能力。\end{analysis}

\q √。\begin{analysis}教材 9.2.3（印刷页 150）指出：脑智能是一种个体智能，而群智能是一种社会智能，或者说是系统智能；微观生理层次上低级的神经元群智能形成了宏观心理层次上高级的脑智能。\end{analysis}

\q ×。改正：符号智能以符号形式的知识和信息为基础，主要通过逻辑推理、运用知识进行问题求解；以数值数据为基础、通过数值计算和算法求解的是计算智能。\begin{analysis}教材 9.2.4（印刷页 150—151）对符号智能和计算智能作了明确区分，本题把两者的基础与求解方式弄反了。\end{analysis}

\q √。\begin{analysis}1956 年，麦卡锡等人在达特茅斯学院会议上提议正式采用“人工智能”这一术语，从而诞生了研究如何用机器模拟人类智能的新兴学科，该会议正式确立了人工智能的研究领域（教材 9.3，印刷页 151）。\end{analysis}

\q ×。改正：图灵测试通过只能说明机器的行为表现与人类相似，不能证明机器具有意识、情感等主观体验。\begin{analysis}图灵测试是一种行为主义意义上的判定方法，它只考察外部对话表现，不考察内部意识（教材第 1 章 1.3.4 仅提及该试验方法；完整含义按通用教材补充）。教材 9.7.2 也指出机器“在意识、推理等方面不可能超越人类”，不能由测试通过推出机器具有意识和情感。\end{analysis}

\q √。\begin{analysis}教材 9.6.14（印刷页 159—160）指出，“深度学习”是指多层的人工神经网络和训练的方法。深度学习属于机器学习的一个分支（教材未系统展开三者层次，按通用教材补充），但神经网络本身不等于深度学习，深度学习强调“多层”和“训练方法”。\end{analysis}

\answerpart{四、配对题（6 对，每对 1 分，共 6 分）}

\q （1）—C\quad （2）—E\quad （3）—A\quad （4）—B\quad （5）—D\quad （6）—F。
\begin{analysis}麦卡洛克与匹茨吸取三种资源后提出人工神经元模型，是人工智能孕育期的最早工作（教材 9.3）；麦卡锡 1956 年提议采用“人工智能”术语（教材 9.3）；费根鲍姆 1977 年提出知识工程（教材 9.3）；布鲁克斯研制六足行走机器人，是行为模拟（控制进化方法）的代表作（教材 9.5.3）；蚁群算法模拟蚂蚁群体觅食活动过程（教材 9.5.4）；符号智能以符号形式的知识和信息为基础，主要通过逻辑推理求解（教材 9.2.4）。\end{analysis}

\answerpart{五、简答题（3 题，每题 8 分，共 24 分）}

\q 参考答案要点：
（1）智能的来源与定义（2 分）：智能一词来源于拉丁语，字面意思是采集、收集、汇集，并由此进行选择。一般认为，智能是指人类在认识世界和改造世界的活动中，由脑力劳动表现出来的能力，即个体对客观事物进行合理分析、判断及有目的的行动和有效处理周围环境事宜的综合能力（教材 9.1，印刷页 149）。
（2）感知能力（2 分）：通过视觉、听觉、味觉、触觉等感觉器官感知外部世界的能力，是获取外部信息的前提，感知是智能活动的必要条件。
（3）记忆和思维能力（2 分）：记忆用于存储感知到的外部信息以及思维产生的知识；思维用于对记忆的信息进行处理，利用已有知识对信息进行分析、计算、比较、判断、联想及对策。思维分为逻辑思维（抽象思维）、形象思维（直感思维）和顿悟思维三种。
（4）学习能力、自适应能力及行为能力（2 分）：学习能力是通过指导、实践等过程丰富自身知识和技巧的能力；自适应能力是在各种环境下都能保持同等效率的能力；行为能力是把思维决策转化为行动的能力。

\q 参考答案要点：
（1）三种定义视角（3 分）：（1）尼尔逊的定义，即人工智能是关于知识的学科，是怎样表示知识以及怎样获得知识并使用知识的科学；（2）从功能定义，人工智能是智能机器所执行的通常与人类智能有关的功能，如判断、推理、证明、识别、学习和问题求解等思维活动；（3）从实用观点看，人工智能是一门知识工程学，以知识为对象，研究知识的获取、知识的表示方法和知识的使用（教材 9.2，印刷页 150）。
（2）一般定义（2 分）：人工智能是研究、开发用于模拟、延伸和扩展人的智能的理论、方法、技术及应用系统的一门新的技术科学；它是计算机科学的一个分支，通过了解智能的实质，生产出能以与人类智能相似的方式做出反应的智能机器（教材 9.2.1）。
（3）智能机器及三要素（3 分）：智能机器是能够在各类环境中自主地或交互地执行各种拟人任务的机器，即让机器会“思考”。它至少要具备感觉要素、反应要素和思考要素；能够理解人类语言并同操作者对话，能分析情况、调整动作，在信息不充分和环境迅速变化的条件下完成动作（教材 9.2.2）。注意教材 9.6.15 讲智能机器人时表述为感觉、运动、思考三要素，答题时可按题源分别说明。

\q 参考答案要点（每点 1.5 分，写出五种即可满分，答全六种不扣分）：
（1）心理模拟，符号推演：从人脑的宏观心理层面入手，以智能行为的心理模型为依据，把问题或知识表示成某种逻辑网络，采用符号推演方法模拟人脑的逻辑思维过程，代表成果是基于知识的智能系统、专家系统（教材 9.5.1）。
（2）生理模拟，神经计算：从人脑的生理层面即微观结构和工作机理入手，以智能行为的生理模型为依据，用人工神经网络作为信息和知识的载体，用神经计算实现网络的学习、记忆、联想、识别和推理（教材 9.5.2）。
（3）行为模拟，控制进化：基于“感知—行为”模型，模拟人和动物在与环境交互、控制过程中的智能活动和行为特性，如反应、适应、学习、寻优；典型代表是 MIT 布鲁克斯研制的六足行走机器人（教材 9.5.3）。
（4）群体模拟，仿生计算：模拟生物群落的群体智能行为，如模拟生物种群有性繁殖和自然选择而出现的遗传算法及进化计算，模拟免疫细胞群的免疫计算、模拟蚂蚁群体觅食的蚁群算法、模拟鸟群飞翔的粒群算法等，通过遗传、变异、选择、交叉、克隆等算子实现（教材 9.5.4）。
（5）博采广鉴，自然计算：从生命、生态、系统、社会、数学、物理、化学乃至经济等众多学科和领域寻找启发和灵感，模拟自然智能开展研究（教材 9.5.5）。
（6）原理分析，数学建模：通过对智能本质和原理的分析，直接采用数学方法建立智能行为模型，如用贝叶斯定理处理不确定性信息，建立统计模式识别、统计机器学习和不确定性推理；用距离、空间、函数、变换等概念开发几何分类、支持向量机等（教材 9.5.6）。

\answerpart{六、综合题（1 题，共 12 分）}

\q 参考答案要点：
（1）智能类型判断（2 分）：属于计算智能。该系统以大量图像数据（数值数据）为基础，用多层神经网络和训练算法求解，符合教材 9.2.4 对计算智能的界定（以数值数据为基础，主要经数值计算、运用算法进行问题求解）；符号智能以符号形式的知识和信息为基础、主要经逻辑推理求解，与该系统不同。
（2）涉及的八大主要内容（3 分，答出三类即可）：（1）感知与交流——通过图像感知病虫害叶片特征，属于机器感知（计算机视觉）；（2）学习与发现——从标注数据中学习识别规律；（3）知识与推理——把识别结果与防治知识结合，给出防治建议；（4）应用与工程——系统直接面向农业生产应用。（教材 9.4，印刷页 152—154）
（3）三者关系（3 分）：人工智能是最大范畴，机器学习是人工智能的分支，研究让机器从数据与经验中获取知识、发现规律；深度学习是机器学习的分支，其模型是多层人工神经网络及训练方法（教材 9.6.14 给出深度学习定义；层次关系教材未展开，按通用教材补充）。该系统用多层神经网络训练，属于深度学习，也属于机器学习和人工智能。
（4）宣传不成立（2 分）：强人工智能要求具备通用智能乃至意识，而该系统只在病虫害识别这一特定任务上表现良好，属于弱人工智能（专用人工智能）；图灵测试通过也只说明行为表现相似，不能证明其具有意识。教材 9.7.2（印刷页 162—163）明确指出，机器人在工作强度、运算速度和记忆功能方面可以超越人类，但在意识、推理等方面不可能超越人类，无法复制创造力、想象力与直觉。
（5）技术背景（2 分）：该系统处于人工智能综合集成期之后、以深度学习和数据驱动为特征的发展阶段。教材 9.6.14 记载 2016 年 AlphaGo 的主要工作原理就是深度学习；教材 9.7.1 的大事记反映了从 1956 年学科诞生、1977 年知识工程、1997 年“深蓝”到 2016 年 AlphaGo 的发展脉络，当前人工智能研究与实际应用结合日益紧密、受应用驱动越来越明显（教材 9.6 末“5 条主线”）。

\end{document}
